\chapter*{\FTMlangFormulaSymbolsName} \label{Formelzeichen}
\addcontentsline{toc}{chapter}{\FTMlangFormulaSymbolsName}
%
%
%
%
%
\begin{flushleft}%
	\begingroup%
	\renewcommand\arraystretch{1.8}%
	\setlength\tabcolsep{0pt}%
	% First Column in math mode!
	\begin{longtable}[l]{>{\fontsize{12pt}{12pt}\selectfont$}M{3.5cm}<{$} >{\fontsize{10pt}{10pt}\selectfont}M{2.3cm}<{} >{\fontsize{10pt}{10pt}\selectfont}m{10cm}<{}}%
		% Header
		\text{\fontsize{10pt}{10pt}\selectfont \FTMlangFormulaSymbolsName} & \text{\FTMlangFormulaSymbolsUnitName} & \FTMlangFormulaSymbolsDescriptionName \\%
		\midrule%
		%% =========================
		%% ----- Insert here -------
		%% =========================
		B_{\mathrm{Fzg}} & \si{\metre} & Breite des Fahrzeuges\\%
		B_{\mathrm{Rad}} & \si{\metre} & Reifenbreite (Breite der Lauffläche)\\%
		D_{\mathrm{Rad}} & \si{\metre} & Raddurchmesser\\%
		H_{\mathrm{Fzg}} & \si{\metre} & Höhe des Fahrzeuges\\%
		R_{\mathrm{Fzg}} & \si{\metre} & Radstand des Fahrzeuges\\%
		S_{\mathrm{Fzg}} & \si{\metre} & Spur des Fahrzeuges\\%
		a_i	& -	& Gewichtungsfaktoren für die Kostenfunktion $J$\\%
		\vc_{\mathrm{DGL}} & - & Gleichungsnebenbedingungen, die aus der Diskretisierung der\newline Zustandsdifferentialgleichung resultieren\\%
		H & - & Hamilton"=Funktion\\%
		E_\mathrm{pot}	& \si{\joule} & Potenzielle Energie\\%
		E_m &  \si{\joule} & Spezifische Energie\\%
		%% =========================
		%% ----- Insert here -------
		%% =========================
	\end{longtable}%
	\endgroup%
\end{flushleft}%
